
\begin{abstract}
This report updates earlier surveys of space-trusted energetic-particle and radiation-monitoring instruments from the SWIMMR Project. It is written as an engineering reference. The review begins with the mature low-Earth-orbit (LEO), medium-Earth-orbit (MEO) and geostationary-orbit (GEO) markets, then considers highly elliptical, interplanetary and emerging cislunar applications. Instruments are compared using factual parameters such as detector technology, particle species, energy coverage, mass, power, volume, field of view, interfaces and flight heritage. Missing values are marked \ND. The final section identifies technology gaps and emerging opportunities. 
\end{abstract}




\section{Purpose and Scope}
The purpose of this survey is to provide engineers and mission designers with an updated, structured comparison of radiation monitors that have flown, are in development, or have recently been proposed. The review consolidates information from the SWIMMR S1 survey, the Kellett trusted-detector summaries, the 2020 radiation-monitoring presentation, and post-2023 literature and mission information. Information contained is limited to that which is publicly available in either peer review papers and organisational websites.

The document focuses on energetic electrons, protons, alpha particles and heavier ions; absorbed dose, dose rate, linear energy transfer (LET), charging and single-event-effect proxies; compact instruments suitable for hosted payloads and small spacecraft; and mission heritage from LEO to the `new market' of lunar and cislunar space.

Energy coverage alone does not establish suitability. Geometric factor, shielding, saturation, contamination, dead time, angular response, calibration and data products remain decisive.

%%%%% why lunar
\paragraph{Scope Extension: Emerging Exploration Market}

In addition to well-established radiation monitors developed for Low Earth Orbit (LEO), Medium Earth Orbit (MEO) and Geostationary Earth Orbit (GEO) missions, this review also includes instruments designed for lunar surface, lunar orbital and cislunar applications. This reflects the rapid expansion of exploration activities beyond Earth orbit, driven by international (China) and commercial programmes such as Artemis, Gateway, Chang'E, Commercial Lunar Payload Services (CLPS), and an increasing number of robotic and commercial missions. 

%The number of planned robotic and human missions operating beyond low Earth orbit is expected to increase substantially over the coming decade.

The coming decade is expected to see more lunar and cislunar missions than at any previous period in the Space Age. This rapid expansion is creating a substantial new market for compact radiation-monitoring instrumentation and is now one of the principal drivers of innovation in detector technology. \cite{NASA_LunarStrategy,CSIS_Cislunar2024}


This is creating new requirements for compact, low-mass and low-power radiation monitors capable of operating in deep-space environments. As a consequence, a significant proportion of current detector innovation is now occurring within these emerging exploration programmes.

Although many of these instruments are being developed specifically for lunar and cislunar missions, they frequently represent the leading edge of radiation-monitor technology. Recent developments in hybrid silicon--scintillator detectors, Timepix-based imagers, compact spectrometers and multi-particle sensor systems have largely been driven by these new mission requirements. As these technologies mature, many are expected to transition into more traditional operational applications in LEO, MEO and GEO.

Accordingly, the inclusion of lunar and cislunar instruments in this review is intended not only to document the current state of the art for exploration missions, but also to identify technologies that are likely to influence the next generation of operational radiation monitors across all orbital regimes.
